% SPDX-License-Identifier: LPPL-1.3c
% Copyright (C) 2026 Christophe Jorssen
% Author and Current Maintainer: Christophe Jorssen <christophe.jorssen@gmail.com>
% LPPL maintenance status: maintained
% This file is part of luafemm. See LICENSE and LICENSES.md for its terms.
% Original tutorial following PGF's progressive, task-led presentation.
\clearpage
\section{Tutorial: A synchronous machine for Kamil}
\label{sec:tutorial-machines}

Kamil is preparing a lesson on synchronous machines. His first drawing has a
circular rotor, an iron stator and two small rectangles for the sides of a
coil. He can explain which way the current flows. What he cannot yet show is
why the field around the air gap should resemble the sine wave on the next
page of his notes. Two rectangles do not look particularly sinusoidal.

He decides to start with those two rectangles, calculate the field, and let
the graph tell him what needs changing. We shall follow him from a concentrated
winding to regularly spaced slots, then to weighted turns and deliberately
unequal slot spacing. Every complete source below is distributed with the
package and executed when this manual is built.

\subsection{Deciding what the drawing represents}

The drawing is a transverse section of a cylindrical machine. The nominal
air-gap surfaces are concentric circles, with no salient poles. Current flows
parallel to the machine axis, out of or into the page. The model has a constant
axial depth. This is a planar field calculation; it is not an axial-flux
machine with a gap between two disks.

Kamil first excites only the stator. The passive iron rotor remains essential:
it changes the magnetic circuit even though its current is zero. Each figure
represents a magnetostatic state. A conducting material does not acquire an
induced current simply because its conductivity is nonzero. End windings,
eddy currents and time-dependent motion are outside this example.

For these comparisons he chooses a linear teaching material with relative
permeability 2000. This avoids mixing winding-distribution effects with
saturation. Later he can replace its identifier with |pure-iron| or another
catalogue material and solve the nonlinear problem.

\subsection{A drawing that is inexpensive to change}

Kamil remembers Maya's first decision: use a modest global mesh while composing
the drawing. Here the global spacing is 4 mm, with 0.7 mm in an air-only band
inside the 2 mm gap. These are working settings, not certified accuracy bounds.
The 0.03 mm curve tolerance controls the polygonal representation of the
circular material interfaces independently of the mesh spacing.

He defines reusable TikZ styles. The component keys use model millimetres,
so |radius=28| is a physical radius; |scale=.65| on the picture only adjusts
the printed diagram. The physical geometry does not grow when he changes its
labels or selects a larger font.

\begin{codeexample}[code only]
\tikzset{
  femm/materials/teaching-iron/.style={mu r=2000},
  kamil rotor/.style={femm/rotor={
    radius=28,core material=teaching-iron}},
  kamil stator/.style={femm/stator={
    bore radius=30,outer radius=45,
    core material=teaching-iron,
    slots={count=2,start angle=90,
      depth=8,neck depth=.8,
      opening width=1,width=3.5,bottom width=4.5},
    coil={slots={1,2},ampere turns=200}}}
}
\end{codeexample}

A slot is a cavity, not a current source. The |coil| record associates two
conductor sections: slot 1 carries $+200$ A\,turns and slot 2 carries
$-200$ A\,turns. Positive current points out of the page. The conductors need
not have equal areas: their current densities are calculated from their
individual captured areas.

\subsection{Two nodes, a gap, and one circular measurement}

The order is worth writing down. Kamil declares the two nodes, adds the gap
refinement, registers the measurement, and finally asks for field lines.
The field request starts the solve. Afterwards he may plot as many components
of the registered profile as he likes, but he cannot add another physical
part to that solved model.

\begin{codeexample}[code only]
\node[kamil stator] (S) {};
\node[kamil rotor] (R) {};
\pic[femm/air gap={rotor=R,stator=S,mesh size=.7}]
  {femm air gap};
\pic[femm/gap profile={rotor=R,stator=S,
  name=gap,samples=721}] {femm gap profile};
\pic[femm/lines=19] {femm field};
\end{codeexample}

The gap profile lies halfway between the nominal surfaces. It samples an
exact circle, uniformly in mechanical angle. Its red drawing is only an
annotation; it does not impose a boundary condition.

The component Kamil wants is $B_r$, positive away from the common centre.
Plotting $|\mathbf B|$ would conceal the change of sign between north and
south poles. |Bn| is also a different convention: the left normal of a
counterclockwise circular path points inward, whereas $B_r$ points outward.

\begin{codeexample}[code only]
\begin{axis}[xlabel={Mechanical angle (degrees)},ylabel={$B_r$ (T)}]
  \femmplot[profile=gap,component=Br,abscissa=angle]
\end{axis}
\end{codeexample}

Here is his complete first document, including the calculated field, the
angular scan and its harmonic spectrum. Its profile is called |kamil-single|
to keep it distinct from the following examples.
\luafemmexample{tutorial-machine-single}

\subsection{More slots do not specify a winding}

Kamil now puts 24 slots around the stator. Their centres start at
$7.5^\circ$ and advance in steps of $15^\circ$. He keeps the opening and body
widths unchanged. With |distribution=uniform|, conductor sections in one
half of the machine carry equal positive ampere-turns, and those in the
opposite half carry equal negative ampere-turns. The selection is based on
the sign of $\sin(\theta)$ for |axis=0,pole pairs=1|.

\begin{codeexample}[code only]
slots={count=24,start angle=7.5,
  depth=8,neck depth=.8,opening width=1,
  width=3.5,bottom width=4.5},
winding={distribution=uniform,ampere turns=200}
\end{codeexample}

The total positive excitation remains 200 A\,turns, as in the concentrated
coil. It is not 200 A\,turns in every slot. This normalization makes the
comparison meaningful, although the fundamental amplitudes need not be equal.
The distributed source produces a different magnetomotive-force waveform.
It does not prescribe the resulting flux density.
\luafemmexample{tutorial-machine-uniform}

\subsection{Changing the turns before moving the slots}

Kamil tries |distribution=sine| with exactly the same slots. For each slot
centre $\theta_i$, the source weight is
\[
 w_i=\sin\bigl(p(\theta_i-\theta_0)\bigr),\qquad
 (NI)_i=(NI)_+\frac{w_i}{\sum_j\max(w_j,0)}.
\]
Here $p$ is |pole pairs| and $\theta_0$ is |axis|. The distribution must be
balanced: $\sum_i w_i=0$ within numerical tolerance. The package rejects an
unbalanced set instead of inventing a return conductor. Multiple layers
share a slot's assigned ampere-turns equally.

In the usual thin-gap, highly permeable-iron approximation, the spatial
variation of magnetomotive force is related to the current distribution.
A sinusoidal current sheet therefore motivates this construction. The FEM
calculation still includes the finite permeability and the actual slots.
This is why the graph need not be a perfect sinusoid. For the continuous
current-sheet model, see the synchronous-machine example in
\href{https://web.mit.edu/6.013_book/www/chapter11/11.7.html}{Haus and Melcher,
\emph{Electromagnetic Fields and Energy}, Section 11.7}.

\begin{codeexample}[code only]
winding={distribution=sine,ampere turns=200,
  pole pairs=1,axis=0}
\end{codeexample}

These weights describe homogenised ampere-turns and can imply noninteger
turn counts at a chosen current. For a manufacturable winding, Kamil would
choose integer turns with explicit |coil| records, then recalculate the
harmonics. No turn rounding is hidden inside |distribution=sine|.
\luafemmexample{tutorial-machine-sine}

\subsection{Keeping equal turns and choosing the positions}

Kamil also wants a construction with identical coil sides. He retains twelve
positive and twelve negative sections, but changes their positions. In the
positive half, he places slot centres at the midpoint quantiles of a density
proportional to $\sin\theta$:
\[
 \theta_i=\arccos\left(1-2\frac{i-1/2}{12}\right),
 \qquad i=1,\ldots,12.
\]
The remaining positions are $\theta_i+180^\circ$. Pairing each slot with its
opposite gives equal positive and negative ampere-turns. The slots are denser
near $90^\circ$ and $270^\circ$, where the continuous source is strongest.

The source below contains the resulting explicit angles. The geometric
checks still apply: adding slots indefinitely would eventually leave too
little room for the teeth. The construction approximates a distribution;
it is not an optimizer and does not guarantee cancellation of particular
harmonics. Its spectrum is the test of the idea.
\luafemmexample{tutorial-machine-positioned}

\subsection{Asking how close is close enough}

The spectrum reports $B_n/B_1$ through order 15. For a component expanded as
$B(\theta)=a_0+\sum_n(a_n\cos n\theta+b_n\sin n\theta)$, the reported amplitude
is $B_n=\sqrt{a_n^2+b_n^2}$. The spatial distortion shown below each example is
\[
 \mathrm{THD}_{15}=\frac{\sqrt{\sum_{n=2}^{15}B_n^2}}{B_1}.
\]
It excludes DC and harmonics above order 15. In particular, a small value
does not imply that slot-scale ripple at higher orders has disappeared.
The output is a ratio; multiply by 100 for percent.

Before using the final figures, Kamil repeats the calculations with global
spacing 2 mm, gap spacing 0.35 mm, and a smaller curve tolerance. He compares
fundamental amplitude, phase, and selected harmonics. He then checks a larger
outer air domain to assess the imposed outer boundary. Increasing |samples|
only improves the scan and Fourier quadrature; increasing |femm/lines| only
adds drawn contours. Neither refines the finite-element solution.

He can finally replace |teaching-iron| with a nonlinear material, change the
stator to a square outer shape, or excite a slotted rotor. Each change is a
new physical comparison. Styles let him keep the same labels, measurement
commands and PGFPlots presentation throughout.
